\section{The Second Week: The Prisoner and the Immaculate Virgin}\label{sec:week-2}

On the Second Sunday of Advent, the Roman station is the Basilica of the
Holy Cross in Jerusalem (\textit{Statio ad S.\ Crucem in Jerusalem}).
The Mass opens with Isaiah's call to the holy city: \textit{Populus
Sion}, ``People of Sion, behold the Lord shall come to save the
nations'' (adapted from Is~30:19, 30), joined to Psalm~79~[80]:2, ``Give
ear, O thou that rulest Israel: thou that leadest Joseph like a sheep.''
In the Gradual (\textit{Ex Sion}, Ps~49~[50]:2--3, 5), the choir sings,
``God shall come manifestly: our God shall come, and shall not keep
silence.'' The Sunday Alleluia (\textit{Laetatus sum}, Ps~121~[122]:1,
dropped on the weekdays of Week~II) turns longing into pilgrimage---``I
rejoiced at the things that were said to me: We shall go into the house
of the Lord''---and the Offertory (\textit{Deus, tu conversus},
Ps~84~[85]:7--8) pleads, ``Thou wilt turn, O God, and bring us to life:
and thy people shall rejoice in thee.'' At Communion (\textit{Jerusalem,
surge}, Bar~5:5; 4:36), Baruch bids Jerusalem rise, stand on high, look
toward the east, and ``behold the joy that cometh to thee from God.''

Where the Collect of the First Sunday cried directly to the Son to rouse
his own power, the Collect of the Second Sunday addresses the Father and
asks him to awaken our hearts:

\begin{massbox}{Collect of the Second Sunday of Advent}
\textbf{Latin.} Excita, Domine, corda nostra ad praeparandas Unigeniti
tui vias: ut, per ejus adventum, purificatis tibi mentibus servire
mereamur: Qui tecum vivit et regnat in unitate.

\textbf{English.} Stir up our hearts, O Lord, to prepare the ways of
Thine only-begotten Son; that through His coming we may attain to serve
Thee with purified minds. Who liveth and reigneth with Thee.
\end{massbox}

Even our preparation for Christ is his Father's gift: we beg God to stir
up our sluggish hearts so that our minds may be pure enough to serve his
only-begotten Son. In the Postcommunion (\textit{Repleti cibo
spiritualis alimoniae}), the Church asks that partaking of the Eucharist
may teach us ``to contemn earthly and love heavenly things.''

\subsection{Patience, Scripture, and the Root of Jesse}

In the Epistle (Rom~15:4--13), St.~Paul states why the Church keeps
reading the ancient prophets through Advent:

\begin{quote}
For what things soever were written were written for our learning:
that, through patience and the comfort of the scriptures, we might have
hope.
\end{quote}

From that foundation Paul urges concord in the Christian assembly:
``Wherefore, receive one another, as Christ also hath received you, unto
the honour of God'' (Rom~15:7). Christ Jesus became ``minister of the
circumcision for the truth of God, to confirm the promises made unto the
fathers'' (Rom~15:8), and by the same mercy gathered the Gentiles to
praise God alongside Israel. Paul weaves four Old Testament texts
together, closing with Isaiah~11:10: ``There shall be a root of Jesse;
and he that shall rise up to rule the Gentiles, in him the Gentiles
shall hope'' (Rom~15:12).

On this same morning, Matins opens Isaiah at chapter~11, beginning with
Isaiah~11:1: ``And there shall come forth a rod out of the root of
Jesse, and a flower shall rise up out of his root.'' Though the 1960
rubrics limit Sunday Matins to Isaiah and St.~Gregory's Gospel homily
(RGBR~n.~220), the passage the Roman Breviary long read in the second
nocturn on this Sunday remains a classic guide to the verse. Writing at
Bethlehem, St.~Jerome explains the rod and the flower
(\textit{Commentary on Isaiah}, Book~IV, ch.~11):

\begin{quote}
We, however, understand that the rod out of the root of Jesse
signifieth the holy Virgin Mary. She was a clean stem that had as yet
put forth no shoot; as we have read above: ``Behold, the Virgin shall
conceive and bear a son.'' (Isa.\ vii.\ 14.) And the flower we believe to
mean the Lord our Redeemer, Who hath elsewhere compared Himself to a
flower; ``I am a flower of the plain, and a lily of the valleys.''
(Cant.\ ii.\ 1.)
\end{quote}

\subsection{John the Baptist in Herod's Prison}

The Gospel of the Second Sunday (Matt~11:2--10) opens with an unsettling
question from a dungeon: ``Now when John had heard in prison the works
of Christ: sending two of his disciples he said to him: Art thou he that
art to come, or look we for another?''

The Fathers stopped over those words. John leaped in Elizabeth's womb at
Mary's voice, saw the Spirit descend at the Jordan, and pointed to Jesus
as the Lamb of God. Why would this same prophet send messengers from
Herod's prison to ask whether Jesus is the Coming One?

St.~John Chrysostom and Pope St.~Gregory the Great both defend John's
unshaken faith while opening two distinct depths in the scene. In his
thirty-sixth homily on Matthew, Chrysostom rules out wavering belief and
fear alike; John had rebuked a tyrant to his face and stood ready for
death: ``Whence it is evident, that neither did he send as being himself
in doubt, nor did he ask in ignorance.'' Knowing his execution was near
and seeing his disciples still clinging to him with jealousy toward
Christ, John sent two of them straight to Jesus so they could see the
works foretold in Isaiah~35:5--6 for themselves---``For the testimony of
deeds is stronger than the testimony of words'' (Chrysostom, in the
\textit{Catena Aurea} on Matt~11:5). Jesus points solely to what is
taking place before them: ``The blind see, the lame walk, the lepers are
cleansed, the deaf hear, the dead rise again, the poor have the gospel
preached to them'' (Matt~11:5).

Once the messengers depart, Jesus turns to the crowds lest anyone think
chains have weakened John: ``What went you out into the desert to see? a
reed shaken with the wind?'' (Matt~11:7). In his thirty-seventh homily
on Matthew, Chrysostom shows Christ calling the people's own memory to
testify to John's steadiness:

\begin{quote}
So much people and so many cities would not have poured themselves out
with so great zeal towards the wilderness and the river Jordan at that
time, had ye not expected to see some great and marvellous one, one
firmer than any rock. Yea, it was not ``a reed'' surely, that ``ye went out
to see shaken by the wind:'' for the flexible and such as are lightly
brought round, and now say one thing, now another, and stand firm in
nothing, are most like that.
\end{quote}

Pope St.~Gregory the Great, whose sixth homily on the Gospels is read
at Matins on the Second Sunday, attends to the order of time (in the
\textit{Catena Aurea} on Matt~11:3). Having already proclaimed Christ at
the Jordan, John asks from prison ``not that he doubted that this was
the Redeemer of the world, but he asks that he may know whether He who
in His own person had come into the world, would in His own person
descend also to the world below''---whether John, going before him in
death, should announce Christ's coming to the righteous awaiting
deliverance below. Gregory also explains why Jesus calls John ``more
than a prophet'' (Matt~11:9; \textit{Catena Aurea}): ``The office of a
prophet is to foretel things to come, not to shew them present. John
therefore is more than a prophet, because Him whom he had foretold by
going before Him, the same he shewed as present by pointing Him out.''

Yet the sharpest edge of Gregory's Sunday homily falls on Christ's
closing beatitude: ``And blessed is he that shall not be scandalized in
me'' (Matt~11:6). Men who admire miracles stumble when the Author of
life hangs upon a cross. ``For the more humbling God hath undergone for
man's sake,'' Gregory preaches (\textit{Homily} 6 on the Gospels), ``the
more worthy is He that man should worship Him.'' Putting Christ's
warning plainly, Gregory continues:

\begin{quote}
Now what is this, but a plain mention of that time, when He afterwards
humbled Himself, becoming obedient unto death, even the death of the
Cross? It is as if He said: ``I indeed do wonderful works, but the day
will come when I shall not refuse to suffer shame and evil treatment.
Take heed then, ye who now worship Me for the works' sake, that when I
come to die ye despise Me not for My death's sake.''
\end{quote}

\subsection{The Feast of the Immaculate Conception (8 December)}

Every Advent includes 8~December, the Feast of the Immaculate Conception
of the Blessed Virgin Mary. In five of the seven calendar shapes it
falls during the Second Week (and on Friday or Saturday of the First
Week when Advent opens on 2~December or~3). Whenever Advent begins on
1~December, 8~December lands directly on the Second Sunday of Advent.
Under the 1962 rubrics, the Immaculate Conception is a first-class feast
standing above the Sundays of Advent in the Table of Precedence
(RG~nn.~15, 91, row~4 against row~6). On that Sunday, therefore, the
Marian feast takes precedence and the Second Sunday receives a
privileged commemoration at Mass and in the Office (RG~nn.~109~a,
111~a). When 8~December falls on an Advent weekday, the weekday likewise
receives a privileged commemoration (RG~n.~109~e). In universal canon
law, 8~December is a holy day of obligation (1983~CIC c.~1246~\S{}1; for
diocesan schedules and transfers, see Section~\ref{sec:penance-law}).

On 8~December, 1854, in the apostolic constitution \textit{Ineffabilis
Deus}, Blessed Pope Pius~IX solemnly defined the dogma: by the authority
of our Lord Jesus Christ, of the blessed Apostles Peter and Paul, and
his own, he declared, pronounced, and defined that the doctrine holding
the Blessed Virgin Mary, in the first instant of her conception, by a
singular grace and privilege of almighty God, in view of the merits of
Jesus Christ, Saviour of the human race, to have been preserved immune
from all stain of original sin, has been revealed by God, and therefore
is to be firmly and constantly believed by all the faithful:

\begin{quote}
\textit{\dots auctoritate Domini Nostri Jesu Christi, beatorum
Apostolorum Petri, et Pauli ac Nostra declaramus, pronunciamus et
definimus, doctrinam, quae tenet, beatissimam Virginem Mariam in primo
instanti suae Conceptionis fuisse singulari omnipotentis Dei gratia et
privilegio, intuitu meritorum Christi Jesu Salvatoris humani generis, ab
omni originalis culpae labe praeservatam immunem, esse a Deo revelatam,
atque idcirco ab omnibus fidelibus firmiter constanterque credendam.}
\end{quote}

On 8~December, the priest vests in white, the organ sounds in festal joy
(Sacred Congregation of Rites, \textit{Instructio de musica sacra},
n.~83~a), and \textit{Gloria in excelsis} breaks Advent's silence. The
Introit \textit{Gaudens gaudebo} sings with Isaiah~61:10, ``I will
greatly rejoice in the Lord, and my soul shall be joyful in my God.''
The Lesson (Prov~8:22--35) contemplates Eternal Wisdom before the
creation of the world, and the Gospel recalls Gabriel's greeting at
Nazareth: ``Hail, full of grace, the Lord is with thee: blessed art thou
among women'' (Luke~1:28). In the Collect, the Church traces Mary's
preservation directly to her Son's foreseen Cross:

\begin{massbox}{Collect of the Immaculate Conception (8 December)}
\textbf{Latin.} Deus, qui per immaculatam Virginis Conceptionem dignum
Filio tuo habitaculum praeparasti: quaesumus; ut, qui ex morte eiusdem
Filii tui praevisa eam ab omni labe praeservasti, nos quoque mundos eius
intercessione ad te pervenire concedas. Per eundem Dominum.

\textbf{English.} O God, Who, by the Immaculate Conception of the Virgin,
didst prepare a worthy dwelling for Thy Son, we beseech Thee that Thou,
Who, by the death, foreseen by Thee, of the same Thy Son, didst preserve
her from all stain, wilt grant us also, by her intercession, to come to
Thee pure in heart. Through the same.
\end{massbox}

The Roman station of the Second Sunday---Holy Cross in Jerusalem---and
the feast of 8~December meet in that phrase \textit{ex morte eiusdem
Filii tui praevisa}: the redeeming death of Christ reached backward in
grace to prepare a stainless dwelling for his own Incarnation.

\subsection{Patience and Hope Through the Second Week}

Across Monday through Saturday, ordinary weekdays without a feast repeat
the Mass of the Second Sunday without the \textit{Alleluia}
(RGMR~n.~299), while third-class feasts such as St.~Peter Chrysologus
(4~December), St.~Nicholas (6~December), and St.~Ambrose (7~December)
take the Mass of the day with a privileged commemoration of the Advent
feria (RG~nn.~25~b, 109~e). At Matins, Isaiah moves from chapter~11
through chapter~25---through the oracles over Babylon, Moab, Egypt, and
the shaken earth (Is~13, 14, 16, 19, 24) to the hymn of deliverance
(Is~25). On Wednesday morning, Matins opens with Isaiah~16:1,
\textit{Emitte Agnum Domine, Dominatorem terrae}: ``Send forth, O Lord,
the lamb, the ruler of the earth, from Petra of the desert, to the mount
of the daughter of Sion.'' On Saturday, Isaiah~25 promises the feast on
the Lord's mountain where ``He shall cast death down headlong for ever:
and the Lord God shall wipe away tears from every face'' (Is~25:8).

Keep the Second Week at home with the Scriptures and the Fathers beside
your daily prayer (see Appendix~\ref{app:daily-course} for each day's
readings and prayer). Read Isaiah~11--25 alongside Romans~15:4--13 and
Matthew~11:2--10, and spend time with St.~Gregory's sixth homily on the
Gospels or St.~John Chrysostom's thirty-seventh homily on Matthew.
Learn the Collect of the Second Sunday (\textit{Excita, Domine, corda
nostra}) by heart and pray it each morning; on 8~December, honour Our
Lady with Mass, festal thanksgiving at table, and the Collect of the
feast (Appendix~\ref{app:prayers}). Put St.~Paul's call to
concord---``receive one another, as Christ also hath received you''
(Rom~15:7)---into daily practice: restrain the impatient word in your
household, visit or write to someone isolated as John was in prison, and
refuse to be scandalized when God works in hidden or humbling ways.
